% AUTO-GENERATED by scripts/build_topics.py — do not edit.
% Edit topics/bodies/topic_11.tex or topics/preamble.tex instead.

\begin{konspekt}
\kreq{Описание}{файлове и основни атрибути; единно пространство от именувани обекти; структура, физическо представяне, точки на монтиране; видове обекти --- файлове, директории, устройства, информационни канали; основни функции --- достъп, изолация и защита, управление на правата --- \kref{11.1}.}
\kreq{Реализация}{кратък обзор на ext2/ext3 --- \kref{11.2}.}
\kreq{Ускоряване и надеждност}{кеширане на вход/изхода, отлагане на записа, алгоритъм на асансьора; журнална файлова система, RAID1 и RAID5 --- \kref{11.3}.}
\kreq{POSIX функции}{\texttt{open}, \texttt{close}, \texttt{read}, \texttt{write}, \texttt{lseek}, включително частичен вход/изход --- \kref{11.4}.}
\kreq{Задача 1}{извличане на интервали от двоичен файл --- схема в \kr{11.5.1}, пълна програма в \kref{11.6}.}
\kreq{Задача 2}{сортиране на файл от 8-битови числа --- \kref{11.5.2}.}
\kreq{Задача 3}{прилагане на двоична кръпка --- \kref{11.5.3}.}
\kstatus{изрично се оценява скорост, пропорционална на реално обработваните данни. Задача 1 е дадена като пълна програма; задачи 2 и 3 са схеми, не програми. Първичните източници за този въпрос липсват в корпуса --- главата следва официалната анотация и подлежи на проверка по литературата.}
\kreq{Литература}{[42, четвърта глава], [48], [49, глави 10--12].}
\end{konspekt}

\begin{supp}[Статус на източниците]
В наличния корпус липсват първичните материали за този въпрос. Главата следва подробната официална анотация и е сверена с предоставения сравнителен конспект. Преди окончателно академично издание тя трябва да се провери по посочената литература: [42, четвърта глава], [48], [49, глави 10--12].
\end{supp}

\section{Файлове и файлова система}\label{s:11.1}

\begin{defn}[Файл]
Файлът е именуван информационен обект, чието съдържание се пази устойчиво и се достъпва като последователност от байтове или записи.
\end{defn}

Операционната система свързва името на файла с метаданни и с физическите блокове, в които е разположено съдържанието. Типични атрибути са тип, размер, собственик, група, права за достъп и времена за последна промяна на съдържанието, метаданните и достъпа. Наличието на отделно време на създаване зависи от конкретната файлова система.

\begin{defn}[Файлова система]
Файловата система е абстракция и реализация за именуване, организиране, съхраняване, намиране, защита и споделяне на информационни обекти.
\end{defn}

\begin{defn}[Видове обекти и директория]
В Unix-подобните системи обектите образуват единно дърво с корен \texttt{/}. В него участват обикновени файлове, директории, символни връзки, специални файлове за устройства, именувани програмни канали (FIFO) и сокети. Директорията съпоставя име с обект на файловата система; самото съдържание на обикновения файл не съдържа името му.
\end{defn}

\begin{defn}[Монтиране]
Отделна файлова система се включва в дървото чрез \emph{монтиране} върху съществуваща директория --- точка на монтиране. Така различни устройства и реализации се виждат през едно пространство от пътища.
\end{defn}

Изолацията и защитата се основават на идентичността на процеса и метаданните на обекта. Класическият Unix модел различава права за четене, запис и изпълнение за собственик, група и останали. За директорията правото за изпълнение означава преминаване и търсене по имена, а правото за запис --- промяна на записите в нея.

\section{Физическо представяне и ext2/ext3}\label{s:11.2}

Дискът се дели на блокове. Файловата система пази служебни структури за свободните и заетите блокове, метаданните на файловете и записите в директориите. При ext2 основни елементи са:

\begin{itemize}
\item \emph{superblock} с общи параметри и състояние на файловата система;
\item групи от блокове с битови карти за свободни блокове и inode-и;
\item таблица от inode структури;
\item блокове с данни и директории.
\end{itemize}

\begin{defn}[Inode]
Inode е структура с метаданните на обекта и указатели към блоковете на съдържанието му. Името се пази в запис на директория, който сочи към номер на inode.
\end{defn}

Малките файлове се адресират чрез преки указатели, а по-големите --- и чрез единично, двойно и тройно косвени блокове с указатели. Твърдата връзка е допълнителен запис на директория към същия inode; символичната връзка е отделен обект, който съдържа път.

\begin{defn}[Журнална файлова система]
ext3 развива ext2 чрез журнал. Преди критични промени те се записват като транзакция в журнала, което позволява след срив файловата система да възстанови съгласувано състояние без пълно сканиране на всички блокове. Режимът на журнализиране определя дали се журнализират само метаданните или и потребителските данни.
\end{defn}

\section{Ускоряване и надеждност}\label{s:11.3}

\begin{defn}[Кеш, отложен запис и алгоритъм на асансьора]
\begin{description}
\item[Кеширане.] Често използвани блокове и файлови страници се пазят в основната памет. Четенето напред (read-ahead) зарежда вероятно нужни следващи блокове.
\item[Отложен запис.] Промените първо се натрупват като замърсени кеширани страници и по-късно се записват групово. Това ускорява работата, но при срив незаписаните данни могат да се загубят; \texttt{fsync} изисква устойчиво записване на съответните промени.
\item[Алгоритъм на асансьора.] При въртящ диск заявките се подреждат подобно на SCAN: главата обслужва позиции в едната посока, после сменя посоката. Целта е по-малко механично преместване; ползата не е същата при SSD.
\end{description}
\end{defn}

\begin{defn}[RAID1 и RAID5]
RAID използва няколко диска като едно блоково устройство. RAID1 огледално дублира данните и понася отказ на един диск във всяка огледална двойка. RAID5 разпределя данните и паритета между поне три диска и обичайно понася отказ на един диск, но записът изисква обновяване на данни и паритет. RAID подобрява наличността, но не е резервно копие: не пази от изтриване, повреда от софтуер или загуба на целия масив.
\end{defn}

\section{POSIX функции за файлове}\label{s:11.4}

\begin{verbatim}
int     open(const char *path, int flags, ...);
int     close(int fd);
ssize_t read(int fd, void *buf, size_t count);
ssize_t write(int fd, const void *buf, size_t count);
off_t   lseek(int fd, off_t offset, int whence);
\end{verbatim}

\texttt{open} връща файлов дескриптор или $-1$. Флаговете задават например четене, запис, създаване и съкращаване. \texttt{read} връща броя реално прочетени байтове, $0$ при край на файл и $-1$ при грешка. \texttt{write} също може да запише по-малко от поисканото, затова надеждната програма повтаря операцията до обработване на целия буфер. Прекъсване със сигнал може да даде \texttt{EINTR} и да изисква повторение.

\texttt{lseek} променя текущото байтово отместване спрямо началото, текущата позиция или края. Тя не извършва вход/изход и не е приложима към всеки тип дескриптор, например към обикновен програмен канал. Всички резултати трябва да се проверяват, а отворените дескриптори да се затварят и при грешка.

\section{Типови задачи}\label{s:11.5}

\subsection{Извличане на интервали от двоичен файл}\label{s:11.5.1}

Файлът \texttt{f1} съдържа двойки $\langle x_i,y_i\rangle$ от \texttt{uint32\_t}, а \texttt{f2} --- масив от такива числа. За всяка двойка се проверява, че интервалът е в границите, позицията се задава на
\[
x_i\cdot\operatorname{sizeof}(\texttt{uint32\_t}),
\]
и се прехвърлят точно $y_i$ числа към \texttt{f3} на блокове. Може да се използва \texttt{lseek} и \texttt{read} или позиционно четене. Не трябва да се обхождат пропуснатите части на \texttt{f2}; времето е пропорционално на броя двойки и размера на изхода. Проверяват се препълване при умножението, частично четене и преждевременен край на файла.

\subsection{Сортиране на файл от 8-битови числа}\label{s:11.5.2}

Понеже има само 256 възможни стойности, се използва броене:

\begin{enumerate}
\item входът се чете последователно на блокове;
\item за всеки байт се увеличава един от 256 брояча с достатъчно широк тип;
\item в изхода се записват последователно съответният брой нули, единици и т.н. до 255.
\end{enumerate}

Така времето е $O(n+256)=O(n)$, а допълнителната памет е $O(256)$ независимо от размер до 2 GB. Изходът също се записва на блокове, а не байт по байт.

\subsection{Прилагане на двоична кръпка}\label{s:11.5.3}

Първо \texttt{f1} се копира в \texttt{f2}. Всеки запис в \texttt{patch} е тройка
\[
(\texttt{uint16\_t offset},\texttt{uint8\_t old},\texttt{uint8\_t new}).
\]
Тройките се обработват последователно. За всяка се проверява, че отместването
съществува и текущият байт в работното копие \texttt{f2} е \texttt{old}; тогава
той се заменя с \texttt{new}. При несъответствие, непълна тройка или I/O грешка
програмата прекратява по контролиран начин.

\begin{supp}[Приета интерпретация: повторено отместване]
Официалната формулировка не казва какво става, когато две тройки сочат едно и
също отместване. Тук е приета следната интерпретация и тя трябва да се \emph{каже
на глас} при отговор: сравнението е винаги с \emph{текущото} състояние на
работното копие \texttt{f2}, а не с неизменния \texttt{f1}. Следствие: при
тройки $(o,a,b)$ и $(o,b,c)$ и двете се прилагат и байтът става $c$; при $(o,a,b)$
и $(o,a,c)$ втората се проваля, защото байтът вече е $b\ne a$. Другата
интерпретация --- сравнение винаги с \texttt{f1} --- е също защитима, но дава
различно поведение и трябва да бъде обявена преди решаването, а не след това.
\end{supp}

Двоичният формат поставя и въпроси за размера и подредбата на числата по байтове. Ако файловете се обменят между различни платформи, форматът трябва изрично да зададе byte order; не бива сляпо да се разчита на представянето на C структура в паметта.

\section{Пълна програма: извличане на интервали}\label{s:11.6}

Разделът по-горе описва решенията в схема. Схема и работеща програма са различни
неща: разликата е в обработката на грешки, частичния вход/изход и затварянето на
дескрипторите. По-долу задача 1 е написана докрай; другите две следват същата
рамка --- сменя се само тялото на цикъла.

\begin{verbatim}
/* Задача 1. f1 съдържа двойки <x, y> от uint32_t;
   f2 е масив от uint32_t. За всяка двойка в f3 се
   записват y числа от f2, започвайки от индекс x.
   Компилиране: cc -std=c11 -O2 -o extract extract.c   */
#include <errno.h>
#include <fcntl.h>
#include <inttypes.h>
#include <stdio.h>
#include <unistd.h>

#define CHUNK 1024

/* read може да върне по-малко от поисканото;
   EINTR не е грешка, а прекъсване.               */
static ssize_t read_full(int fd, void* buf, size_t n) {
  size_t done = 0;
  while (done < n) {
    ssize_t r = read(fd, (char*)buf + done, n - done);
    if (r < 0) {
      if (errno == EINTR) continue;
      return -1;
    }
    if (r == 0) break;                    /* край на файла */
    done += (size_t)r;
  }
  return (ssize_t)done;
}

/* write също може да запише частично --- това е
   най-често пропусканото място.                 */
static int write_full(int fd, const void* buf, size_t n) {
  size_t done = 0;
  while (done < n) {
    ssize_t w = write(fd, (const char*)buf + done, n - done);
    if (w < 0) {
      if (errno == EINTR) continue;
      return -1;
    }
    done += (size_t)w;
  }
  return 0;
}

int main(int argc, char** argv) {
  if (argc != 4) {
    fprintf(stderr, "употреба: %s f1 f2 f3\n", argv[0]);
    return 2;
  }

  int f1 = open(argv[1], O_RDONLY);
  int f2 = open(argv[2], O_RDONLY);
  int f3 = open(argv[3], O_WRONLY | O_CREAT | O_TRUNC, 0644);
  if (f1 < 0 || f2 < 0 || f3 < 0) {
    perror("open");
    return 1;
  }

  /* колко числа има в f2 */
  off_t size2 = lseek(f2, 0, SEEK_END);
  if (size2 < 0) {
    perror("lseek");
    return 1;
  }
  uint64_t count2 = (uint64_t)size2 / sizeof(uint32_t);

  int status = 0;
  uint32_t pair[2];
  ssize_t got;

  const ssize_t PAIR = (ssize_t)sizeof pair;

  while ((got = read_full(f1, pair, sizeof pair)) == PAIR) {
    uint64_t from = pair[0];
    uint64_t len  = pair[1];

    /* Проверката е написана така, че from + len да не прелее. */
    if (from > count2 || len > count2 - from) {
      fprintf(stderr, "интервал извън f2: [%" PRIu64
              ", %" PRIu64 ")\n", from, from + len);
      status = 1;
      break;
    }

    /* Позиционираме се; пропуснатите части на f2
       изобщо не се четат. */
    if (lseek(f2, (off_t)(from * sizeof(uint32_t)), SEEK_SET) < 0) {
      perror("lseek");
      status = 1;
      break;
    }

    uint32_t buf[CHUNK];
    while (len > 0) {
      size_t take  = len < CHUNK ? (size_t)len : CHUNK;
      size_t bytes = take * sizeof(uint32_t);
      if (read_full(f2, buf, bytes) != (ssize_t)bytes) {
        fprintf(stderr, "непълно четене от f2\n");
        status = 1;
        break;
      }
      if (write_full(f3, buf, bytes) != 0) {
        perror("write");
        status = 1;
        break;
      }
      len -= take;
    }
    if (status != 0) break;
  }

  if (got < 0) {
    perror("read");
    status = 1;
  } else if (status == 0 && got != 0) {
    fprintf(stderr, "непълна двойка в края на f1\n");
    status = 1;
  }

  if (close(f1) < 0 || close(f2) < 0 || close(f3) < 0) {
    /* грешка при close е загуба на данни */
    perror("close");
    status = 1;
  }
  return status;
}
\end{verbatim}

Четирите решения, които отличават работеща програма от схема:

\begin{itemize}
\item \texttt{read\_full} и \texttt{write\_full} --- частичен вход/изход и \texttt{EINTR}
се обработват на едно място, вместо да се повтарят във всеки цикъл;
\item проверката \verb|len > count2 - from| вместо \verb|from + len > count2| ---
втората прелива при големи стойности и пропуска невалиден интервал;
\item \texttt{lseek} вместо четене на пропуснатите части --- времето остава
пропорционално на броя двойки плюс размера на изхода, а не на размера на \texttt{f2};
\item резултатът на \texttt{close} се проверява: при отложен запис грешката се
проявява точно там.
\end{itemize}

\begin{supp}[Обхват]
Другите две задачи (\S11.5.2 и \S11.5.3) са дадени като схеми, не като пълни
програми. За сортирането на байтове това е достатъчно --- схемата от три реда е
целият алгоритъм и рамката по-горе се пренася дословно. За кръпката пълната
програма изисква и разрешаването на въпроса с повторените отмествания, обявен
по-горе.
\end{supp}

\section{Изпитен фокус}\label{s:11.7}

\begin{itemize}
\item Файловата система като единно именувано дърво; обекти, метаданни, права и точки на монтиране.
\item Ролята на superblock, inode, записите в директориите и преките/косвените блокове при ext2; журналът при ext3.
\item Кеш, read-ahead, отложен запис, алгоритъм на асансьора, RAID1 и RAID5; защо RAID не е backup.
\item Семантика и проверка на резултатите на \texttt{open}, \texttt{close}, \texttt{read}, \texttt{write} и \texttt{lseek}, включително частичен вход/изход.
\item Решение на трите типови задачи с време, пропорционално на реално обработваните данни.
\item Пълна програма за задача 1 (\S11.6): цикли за частичен \texttt{read}/\texttt{write}, обработка на \texttt{EINTR}, проверка без препълване, \texttt{lseek} вместо обхождане и проверка на \texttt{close}.
\item Обявяване на приетата интерпретация при повторено отместване в задачата за кръпката, преди решаването.
\end{itemize}
